\section{Materiality y Disclosure: desde cuándo se cuentan los cuatro días hábiles}

El equipo de seguridad puede evaluar el impacto técnico, pero no puede decidir por sí solo todas las obligaciones de divulgación. La \term{materiality} (importancia relevante), en el contexto de la legislación bursátil de Estados Unidos, se centra en si es sustancialmente probable que un inversionista razonable considere esa información importante para su decisión de inversión; no la determinan por sí solos el número de registros, el CVSS, los minutos de interrupción ni la etiqueta del atacante. Además, cada jurisdicción, sector, contrato y regla de notificación a usuarios tiene su propio trigger, por lo que la organización necesita un parallel analysis en lugar de buscar un «umbral universal».

\subsection{Facts, materiality, audience y timing}

El flujo de trabajo de divulgación establece primero hechos confiables; después, las personas adecuadas analizan el impacto operativo, financiero, reputacional, legal y estratégico, y luego se identifican los destinatarios y los plazos. Los destinatarios pueden incluir a los usuarios afectados, socios de negocio, aseguradoras, autoridades policiales, reguladores de privacidad o del sector, y el mercado público. Las distintas notificaciones pueden tener propósitos y niveles de detalle diferentes, pero no pueden contradecirse. Al leer este flujo conviene notar en particular lo siguiente: la technical severity solo responde con qué rapidez deben asignarse recursos de defensa, mientras que la materiality responde si el incidente es importante para destinatarios legales y de decisión específicos; la primera puede detonar la segunda, pero no puede sustituirla automáticamente.

\lecturefigure{08-disclosure-decision.png}{La decisión de divulgación parte de los hechos técnicos, pasa por el análisis de materiality y llega a la elección de destinatarios, contenido y plazos.}{SEC 2023 cybersecurity disclosure final rule; NIST guidance; rediseño conceptual.}

\begin{knowledgebox}{Lectura de la figura: primero se determina, después se inicia el reloj correspondiente}
Facts enumera scope, data, systems, operational impact y confidence; materiality analiza investor/customer/user harm, financial/operational effect y cumulative impact; audience enumera users, regulators, law enforcement y partners; timing hace después la correspondencia con contract, statute, Form 8-K o approved delay. El orden esencial es «tras el descubrimiento, determinar la importancia relevante sin demora injustificada; una vez determinada como relevante, calcular la ventana de presentación correspondiente», y no que todos los incidentes tengan uniformemente cuatro días a partir de su descubrimiento.
\end{knowledgebox}

\begin{importantbox}{Alcance preciso de la regla de cuatro días hábiles de la SEC}
Para las empresas estadounidenses que cotizan en bolsa y a las que aplica, la regla de la SEC exige determinar la importancia relevante sin demora injustificada después de descubrir el incident; si la empresa determina que el incidente es relevante, normalmente debe presentar el Form 8-K Item 1.05 \textbf{dentro de los cuatro días hábiles posteriores a esa determinación de importancia relevante}. El \term{Form 8-K} es el formulario de reporte corriente con el que las empresas cotizadas informan sobre ciertos eventos importantes. La regla incluye un mecanismo de aplazamiento para cuando el Procurador General de Estados Unidos determine que la divulgación inmediata representaría un riesgo sustancial para la seguridad nacional o la seguridad pública. Su aplicación concreta debe evaluarla un counsel calificado.
\end{importantbox}

\begin{warningbox}{No hay que esperar a que «todos los hechos sean cien por ciento seguros»}
La investigación puede durar semanas, pero el plazo de notificación puede ser más corto. El objetivo de la gobernanza no es usar la incertidumbre para aplazar, sino dejar claro lo conocido, lo desconocido, el nivel de confianza y las actualizaciones previstas. Dar cifras inexactas demasiado pronto daña la confianza, y poner en marcha demasiado tarde el materiality workstream también crea riesgo de incumplimiento. Hay que involucrar al disclosure owner desde las primeras etapas del incidente y establecer checkpoint de reevaluación periódica.
\end{warningbox}

\begin{knowledgebox}{Cómo evita RACI la dilución de responsabilidades}
\term{RACI} es la matriz de responsabilidades Responsible, Accountable, Consulted, Informed. Para el technical scoping, el forensics lead puede ser Responsible y el IC Accountable; para las obligaciones legales, el counsel es Responsible y el GC Accountable; para la divulgación de asuntos operativos importantes, el CEO o el disclosure committee es Accountable. RACI no busca que los cuatro tipos de rol carguen la responsabilidad por partes iguales, sino dejar claro quién es el único responsable final de cada deliverable.
\end{knowledgebox}

\subsection{Resumen de la sección}

Disclosure es una decisión basada en evidencia, interfuncional y sujeta a varios plazos. El equipo de seguridad aporta hechos oportunos y confiables, el counsel interpreta las obligaciones, y la alta dirección o el comité asume las decisiones operativas y de divulgación pública; los «cuatro días» comienzan con la determinación de importancia relevante, no son una cuenta regresiva mecánica a partir de la primera alerta.
